\chaptermark{Clôture cardinale bidirectionnelle du continu analytique HMT}
\label{chap:83-19}

Ce chapitre ne propose pas de générateur alternatif du continu. Sa prémisse
causale est le bloc unique déjà construit
\[
 \begin{aligned}
 \mathrm{APP}&\longrightarrow\mathrm{TRIT}\longrightarrow\mathrm{TPK}
 \longrightarrow\text{état enrichi unique}\\
 &\longrightarrow\text{état commun, arbre et mesure généalogique}\\
 &\longrightarrow\text{opérations intrinsèques corrélatives}\\
 &\longrightarrow\text{assemblage naturel, terminalité et limite}
 \longrightarrow\mathcal C_{\rm cont}^{\rm disc}.
 \end{aligned}
\]
Les cinq constructions consubstantielles développées dans le chapitre
précédent ont émergé corrélativement sur ce même état avant la
publication réelle, sans prétendre épuiser toutes ses propriétés. Ce qui suit est une présentation
ultérieure du quotient, de sa mémoire temporelle et de sa puissance ; elle ne sélectionne
rétrospectivement aucune de ces constructions.

\section{La cellule TRIT complète et l'extension temporelle bidirectionnelle}

Le signe temporel ne se sépare pas de l'état qu'il oriente. La cellule locale de
départ est
\begin{equation}
 c=(i,j;R^+,Q^+,R^\times,Q^\times;\tau),
 \qquad \tau\in\{-1,0,+1\},
 \label{p12-83-c19-s1:hmtch:eq-celula-trit-completa}
\end{equation}
où
\[
 R^+=\operatorname{dr}_9(i+j),\qquad
 Q^+=\frac{i+j-R^+}{9},
\]
\[
 R^\times=\operatorname{dr}_9(ij),\qquad
 Q^\times=\frac{ij-R^\times}{9}.
\]
Les résidus constituent la lecture visible et les quotients conservent la
mémoire nécessaire pour inverser le transport. Le trit temporel agit sur
cette cellule entière :
\begin{equation}
\begin{aligned}
 +1&\longmapsto \text{retour, mémoire retenue et passé},\\
 0&\longmapsto \text{seuil d'évaluation et présent},\\
 -1&\longmapsto \text{ouverture bidirectionnelle et futur}.
\end{aligned}
\label{p12-83-c19-s1:hmtch:eq-lector-temporal-trit}
\end{equation}
Ainsi, deux occurrences du même signe ne sont pas le même état si leurs
quotients, leur feuillet, leur retenue ou leur registre généalogique diffèrent.

Soit \(X\) une carte compacte de la limite survivante APP--TRIT--TPK et soit
\(F:X\to X\) sa mise à jour complète, continue et surjective. La
surjectivité des
troncatures et la prolongation survivante permettent de former l'extension
naturelle
\begin{equation}
 X_F^{\leftrightarrow}
 :=\left\{(x_k)_{k\in\mathbb Z}\in X^{\mathbb Z}:
 F(x_k)=x_{k+1}\ \text{pour tout }k\in\mathbb Z\right\}.
 \label{p12-83-c19-s1:hmtch:eq-extension-natural-two-way}
\end{equation}
Le déplacement
\[
 \sigma((x_k)_{k\in\mathbb Z})=(x_{k+1})_{k\in\mathbb Z}
\]
ne perd plus la direction inverse.

\begin{theorem}[Extension naturelle two-way de l'état enrichi]
\label{p12-83-c19-s1:hmtch:thm-extension-natural-two-way}
L'espace \(X_F^{\leftrightarrow}\) est compact et

\[
 \sigma:X_F^{\leftrightarrow}\longrightarrow X_F^{\leftrightarrow}
\]
est un homéomorphisme, avec
\[
 \sigma^{-1}((x_k))=(x_{k-1}).
\]
Par conséquent, aller et retour sont des morphismes continus inverses sur l'objet
enrichi. La discontinuité apparente n'apparaît que si l'on oublie d'abord
les coordonnées qui rendent l'histoire inversible.
\end{theorem}

\begin{proof}
Le produit \(X^{\mathbb Z}\) est compact. La condition
\(F(x_k)=x_{k+1}\) est fermée puisque \(F\) est continue ; donc
\(X_F^{\leftrightarrow}\) est compact. Le décalage et le
décalage inverse sont continus dans la topologie produit et préservent
les équations de compatibilité. Leurs compositions sont l'identité.
\end{proof}

L'inversion matérielle de la route vit dans la complétion en groupoïde du
transport. Soit \(C_M=-\operatorname{rev}\) la conjugaison de mémoire. Son
action contravariante s'exprime par le carré de réversion
\[
 F(x_k)=x_{k+1}
 \quad\Longrightarrow\quad
 F(C_Mx_{k+1})=C_Mx_k.
\]
Sur l'extension naturelle, on définit
\begin{equation}
 (\Theta u)_k:=C_M(x_{-k}),
 \qquad u=(x_k)_{k\in\mathbb Z}.
 \label{p12-83-c19-s1:hmtch:eq-definicion-espejo-two-way}
\end{equation}
L'identité de réversion garantit que \(\Theta u\) satisfait de nouveau
les équations de compatibilité. Par conséquent,
\begin{equation}
 \Theta^2=I,
 \qquad
 \Theta\sigma\Theta=\sigma^{-1}.
 \label{p12-83-c19-s1:hmtch:eq-espejo-inversion-two-way}
\end{equation}
Cette identité n'identifie pas le groupe ternaire au signe binaire : le
TRIT oriente retour, seuil et ouverture, tandis que l'involution \(C_M\)
met en œuvre la conjugaison des deux orientations temporelles.

Pour \(n\geq0\), notons \(S_n^{\leftrightarrow}\) l'ensemble fini
des fenêtres temporelles admissibles
\((x_{-n},\ldots,x_0,\ldots,x_n)\), chaque état étant en outre tronqué à la
profondeur généalogique \(n\). Les restrictions simultanées de temps et de
profondeur forment un
système projectif et
\begin{equation}
\mathscr O_{\leftrightarrow}^{\rm HMT}
 :=\varprojlim_n S_n^{\leftrightarrow}
 \cong X_F^{\leftrightarrow}.
 \label{p12-83-c19-s1:hmtch:eq-biordinal-genealogico}
\end{equation}
Chaque histoire possède un support temporel de type
\(\omega^*+1+\omega\), mais porte en outre la cellule
\eqref{p12-83-c19-s1:hmtch:eq-celula-trit-completa}, la route et le registre généalogique. L'application
qui envoie une orbite sur toutes ses fenêtres centrées est un homéomorphisme et
entrelace \(\sigma\) avec la succession généalogique. C'est le biordinal
généalogique : non l'ordinal de von Neumann obtenu après oubli des
étiquettes.

L'objet précédent est une carte compacte de la partie fractionnaire. L'objet
global qui publie toute la droite est l'atlas dénombrable
\begin{equation}
 \mathscr O_{\mathbb R,\leftrightarrow}^{\rm HMT}
 :=\mathbb Z\times\mathscr O_{\leftrightarrow}^{\rm HMT}.
 \label{p12-83-c19-s1:hmtch:eq-biordinal-global}
\end{equation}
La coordonnée \(\mathbb Z\) ne fait que translater la carte ; elle n'efface pas le registre généalogique
bidirectionnel.

\section{Calendrier unitaire, mémoire modulaire et mode stable}

La structure temporelle ne repose pas uniquement sur une métaphore du retour.
Sur
\(\mathcal H_{\rm clk}=\mathbb C^{108}\), soit
\(\zeta=\exp(2\pi i/108)\),
\[
 U_{\rm clk}|j\rangle=|j+1\rangle,
 \qquad
 |\widetilde k\rangle
 =\frac1{\sqrt{108}}\sum_{j=0}^{107}\zeta^{jk}|j\rangle,
\]
et définissons
\begin{equation}
 H_{\rm clk}
 =\hbar_{\rm int}\omega_0
 \sum_{k=0}^{107}k\,|\widetilde k\rangle\langle\widetilde k|,
 \qquad
 \omega_0=\frac{2\pi}{108t_0}.
 \label{p12-83-c19-s1:hmtch:eq-hamiltoniano-reloj}
\end{equation}
Alors
\begin{equation}
 U_{\rm step}
 =\exp\!\left(-\frac{i}{\hbar_{\rm int}}H_{\rm clk}t_0\right)
 =U_{\rm clk},
 \qquad U_{\rm step}^{108}=I,
 \label{p12-83-c19-s1:hmtch:eq-propagador-reloj}
\end{equation}
et, pour \(Z_{\rm clk}|j\rangle=\zeta^j|j\rangle\),
\[
 |\psi_n\rangle:=U_{\rm step}^n|j_0\rangle,
\]
\begin{equation}
 \langle Z_{\rm clk}\rangle_n
 :=\langle\psi_n|Z_{\rm clk}|\psi_n\rangle
 =\zeta^{j_0+n}
 \label{p12-83-c19-s1:hmtch:eq-respuesta-primitiva-108}
\end{equation}
a pour période primitive \(108\).

Dans le secteur modulaire de bord
\(\mathcal H_A=(\mathbb C^3)^{\otimes36}\), soit \(\rho_A>0\) précisément
l'état fidèle caractérisé par la décomposition qui suit et soit
\(\Delta_4^{\rm mod}\) le saut modulaire déclaré. Le générateur qui conserve
la provenance des deux canaux du bord est
\begin{equation}
 K_{\rm HHM}^{(A)}=-\log\rho_A
 =\kappa_A I+\Delta_4^{\rm mod}(90N_{90}+120N_{120}).
 \label{p12-83-c19-s1:hmtch:eq-hamiltoniano-modular-memoria}
\end{equation}
Si
\[
 N=N_{90}+N_{120},
 \qquad
 D=90N_{90}+120N_{120},
\]
alors
\begin{equation}
 N_{90}=4N-\frac{D}{30},
 \qquad
 N_{120}=\frac{D}{30}-3N.
 \label{p12-83-c19-s1:hmtch:eq-reconstruccion-canales-modulares}
\end{equation}
Ainsi, l'observable totale et le hamiltonien modulaire reconstruisent
conjointement les deux mémoires de provenance. Par ailleurs, la dynamique
cylindrique de Perron conserve le mode
\begin{equation}
 V_n=\varphi^n,
 \qquad
 M_n=\varphi^{-2n}=V_n^{-2},
 \qquad
 M_n=\left(\frac{a_n}{a_0}\right)^{-6}\qquad(D=3).
 \label{p12-83-c19-s1:hmtch:eq-modo-estable-memoria}
\end{equation}

\begin{theorem}[Composition finie exacte du calendrier CT108 et de la mémoire
modulaire \(90/120\)]
\label{p12-83-c19-s1:hmtch:thm-composicion-calendario-memoria}
Les opérateurs
\eqref{p12-83-c19-s1:hmtch:eq-hamiltoniano-reloj} et
\eqref{p12-83-c19-s1:hmtch:eq-hamiltoniano-modular-memoria} sont autoadjoints dans leurs
espaces respectifs. Une fois fixée une échelle énergétique \(E_\partial>0\),
l'opérateur du produit
\begin{equation}
 H_{\rm tw}
 :=H_{\rm clk}\otimes I
 +I\otimes E_\partial(K_{\rm HHM}^{(A)}-\kappa_A I)
 \label{p12-83-c19-s1:hmtch:eq-hamiltoniano-two-way-total}
\end{equation}
est autoadjoint et ses deux termes commutent fortement. Pour les états dont
la marginale d'horloge est \(|j_0\rangle\langle j_0|\), l'observable
\(Z_{\rm clk}\otimes I\) conserve sa réponse primitive de période \(108\).
Simultanément, \(K_{\rm HHM}^{(A)}\), \(N_{90}\) et \(N_{120}\) sont
des observables conservées du facteur modulaire. Le couple \((N,D)\), avec
\(N=N_{90}+N_{120}\), reconstruit exactement \((N_{90},N_{120})\) au moyen de
\eqref{p12-83-c19-s1:hmtch:eq-reconstruccion-canales-modulares}.
\end{theorem}

\begin{proof}
Le premier opérateur possède une décomposition spectrale réelle finie, et le
second est le logarithme d'un état fidèle ; tous deux sont autoadjoints. Ils agissent sur
des facteurs tensoriels distincts et, par conséquent, leurs résolutions spectrales
commutent. L'exponentielle de leur somme se factorise. La primitivité de la réponse
provient de la primitivité de \(\zeta\) ; la reconstruction des canaux est
l'inversion de la matrice
\(\left(\begin{smallmatrix}1&1\\90&120\end{smallmatrix}\right)\), de
determinante \(30\).
\end{proof}

\begin{corollary}[Mode stable de la mémoire cylindrique]
\label{p12-83-c19-s1:hmtch:cor-modo-estable-memoria}
Dans le secteur de Perron des routes fermées,
\(M_n=\varphi^{-2n}=V_n^{-2}\) est le mode spectral exact et stable donné
par \eqref{p12-83-c19-s1:hmtch:eq-modo-estable-memoria}.
\end{corollary}

Le calendrier CT108, la mémoire modulaire et le mode de Perron sont trois lecteurs
typés du même antécédent généalogique ; ce ne sont pas le même opérateur. Avec
l'extension naturelle et le miroir, ils forment le registre temporel interne
\[
 \mathcal T_{\rm HMT}
 :=\bigl(X_F^{\leftrightarrow},\sigma,\Theta;
 H_{\rm clk},K_{\rm HHM}^{(A)};M_\bullet\bigr).
\]
Tel est le contenu exact de la clôture temporelle HMT : deux directions
continues inverses, une horloge primitive, une mémoire récupérable et un mode stable.
L'état TPK complet n'est pas remplacé par le curseur \(108\). Une identification
supplémentaire des trois lecteurs à une unique dynamique physique perturbativement
stable exigerait un entrelaceur commun ; cette identification n'est utilisée ni dans
la clôture généalogique ni dans la reconstruction du continu.

\section{Publications analytiques d'histoires HMT}

Soit \(H\) un espace polonais d'histoires : une limite projective fermée de
niveaux finis ou un atlas dénombrable de telles limites. La topologie est
engendrée par les cylindres de préfixe. Ce choix conserve la généalogie et
permet de typer exactement les opérations descriptives admises.

\begin{definition}[Publication HMT analytique]
\label{p12-83-c19-s1:hmtch:def-publicacion-analitica}
Une publication analytique est un ensemble \(A\subseteq\mathbb R\) pour lequel
il existe un borélien
\[
 C\subseteq H\times\mathbb N^{\mathbb N}
\]
et un lecteur borélien \(F:H\to\mathbb R\) tels que
\[
 A=\{F(x):\exists z\in\mathbb N^{\mathbb N},\ (x,z)\in C\}.
\]
Le témoin \(z\) peut coder une route, une prolongation, une mémoire ou
une certification dénombrable. Les lecteurs continus sont des cas particuliers.
\end{definition}

Cette classe contient les cylindres, les prédicats boréliens sur les
histoires, les opérations dénombrables, les images par des lecteurs continus
et les publications qui affirment l'existence d'un témoin réel. Elle ne
contient pas par définition tout l'ensemble des parties \(\mathcal P(\mathbb R)\) : les
réunions non dénombrables, les sélecteurs arbitraires, les quotients non lisses et
la récursion de longueur \(\omega_1\) exigent des données supplémentaires.

\begin{theorem}[Dichotomie parfaite des publications HMT]
\label{p12-83-c19-s1:hmtch:thm-dicotomia-perfecta}
Toute publication HMT analytique \(A\subseteq\mathbb R\) satisfait exactement
l'une des deux alternatives :
\[
 |A|\leq\aleph_0
 \qquad\text{ou}\qquad
 |A|=2^{\aleph_0}.
\]
Dans la seconde alternative, il existe une injection continue
\[
 j:2^{\mathbb N}\hookrightarrow A.
\]
Par conséquent, toute publication analytique non dénombrable contient une copie de
l'espace de Cantor et, en particulier, un sous-ensemble parfait.
\end{theorem}

\begin{proof}
La projection existentielle d'un borélien d'un produit d'espaces polonais
est analytique, et une image borélienne d'un espace borélien standard est de
nouveau analytique. Par conséquent, \(A\) est un sous-ensemble analytique de
\(\mathbb R\).

Pour expliciter l'alternative non dénombrable, représentons \(A\) comme
la projection d'un fermé
\[
 D\subseteq\mathbb R\times\mathbb N^{\mathbb N}.
\]
Pour chaque mot fini \(s\in\mathbb N^{<\mathbb N}\), soit
\[
 A_s:=\{x:\exists z\supseteq s,\ (x,z)\in D\}.
\]
Si \(A_s\cap K\) est non dénombrable pour un certain intervalle fermé \(K\),
l'ensemble de ses points qui ne sont pas des points de condensation est dénombrable, car
\(\mathbb R\) possède une base dénombrable. On choisit deux points de condensation
distincts, deux intervalles fermés disjoints à l'intérieur de \(K\), et pour chaque
intervalle une prolongation stricte de \(s\) dont la projection demeure non
dénombrable. Les deux prolongations n'ont pas besoin d'être incompatibles : la
séparation des branches est garantie par les intervalles. L'identité
\[
 A_s=\bigcup_{n\in\mathbb N}A_{s^\frown n}
\]
garantit que, dans chaque voisinage choisi, une extension conserve la non-
dénombrabilité.

En répétant la procédure, on obtient, pour chaque
\(u\in2^{<\mathbb N}\), un intervalle \(K_u\) et un préfixe \(s_u\) tels que
\[
 K_{u0}\cap K_{u1}=\varnothing,
 \qquad
 K_{ui}\subseteq K_u,
 \qquad
 \operatorname{diam}K_u\leq2^{-|u|},
\]
et \(s_{ui}\) prolonge strictement \(s_u\). Pour
\(\beta\in2^{\mathbb N}\), les intervalles
\(K_{\beta|n}\) déterminent un unique point \(j(\beta)\), et les préfixes
\(s_{\beta|n}\) déterminent un témoin \(z_\beta\). La fermeture de \(D\)
implique
\[
 (j(\beta),z_\beta)\in D,
\]
de sorte que \(j(\beta)\in A\). Les intervalles frères disjoints prouvent
l'injectivité, et le contrôle des diamètres prouve la continuité de \(j\).

Ainsi,
\[
 2^{\aleph_0}\leq |A|\leq|\mathbb R|=2^{\aleph_0},
\]
ce qui clôt la seconde alternative. Si le processus ne part pas d'un ensemble
non dénombrable, on obtient la première.
\end{proof}

\begin{corollary}[Hypothèse du continu analytique HMT]
\label{p12-83-c19-s1:hmtch:cor-ch-analitica-hmt}
Soit
\begin{equation}
 \aleph_{1,\mathrm{HMT}}^{\rm an}
 :=\min\left\{|A|:
 A\text{ est une publication HMT analytique non dénombrable}\right\}.
 \label{p12-83-c19-s1:hmtch:eq-aleph1-analitico-hmt}
\end{equation}
Ce symbole désigne le cardinal initial relatif à la classe des publications
analytiques HMT ; il ne désigne pas l'ordinal extensionnel \(\omega_1\).
Alors
\begin{equation}
 \boxed{
 \aleph_{1,\mathrm{HMT}}^{\rm an}
 =2^{\aleph_0}
 =|\mathbb R|
 =\left|\mathscr O_{\mathbb R,\leftrightarrow}^{\rm HMT}/
 \!\sim_{\mathfrak R}\right|.}
 \label{p12-83-c19-s1:hmtch:eq-ch-analitica-hmt}
\end{equation}
Autrement dit, il n'existe pas de cardinal intermédiaire entre le dénombrable et le continu
au sein des publications analytiques produites par les lecteurs HMT.
\end{corollary}

\begin{proof}
L'ensemble des fenêtres de profondeur finie
\[
 \mathscr O_{<\omega}^{\rm HMT}
 :=\bigcup_{n<\omega}S_n^{\leftrightarrow}
\]
est dénombrable, puisqu'il est une réunion dénombrable d'ensembles finis. La couche complète
se surjecte sur \(\mathbb R\) par
\eqref{p12-83-c18-s2:hmtch:eq-homeomorfismo-two-way-real}, et est contenue dans un produit
dénombrable d'alphabets finis ; elle a donc pour cardinal
\(2^{\aleph_0}\). Le
\cref{p12-83-c19-s1:hmtch:thm-dicotomia-perfecta} exclut toute taille analytique
strictement intermédiaire.
\end{proof}

\begin{corollary}[Dichotomie des fibres]
\label{p12-83-c19-s1:hmtch:cor-dicotomia-fibras}
Soit \(F:H\to Y\) un lecteur borélien entre espaces polonais et soit
\(y\in Y\). Si le domaine des histoires admises est analytique, la fibre
\[
 \{x:F(x)=y\}
\]
est dénombrable ou possède le cardinal \(2^{\aleph_0}\).
\end{corollary}

\begin{proof}
La fibre est l'intersection du domaine analytique avec l'image réciproque
borélienne du singleton \(\{y\}\) ; elle est donc analytique. On applique le
\cref{p12-83-c19-s1:hmtch:thm-dicotomia-perfecta}.
\end{proof}

\section{Comparaison entre la clôture cardinale généalogique et sa réduction ordinale extensionnelle}

Soit \(\operatorname{WO}\subseteq2^{\mathbb N}\) l'ensemble des codes
réels des bons ordres dénombrables, en permettant que le champ soit un
sous-ensemble de \(\mathbb N\). L'application
\[
 \operatorname{ot}:\operatorname{WO}\longrightarrow\omega_1
\]
attribue à chaque code son type d'ordre et est surjective.

\begin{proposition}[Injection ordinale dans la droite]
\label{p12-83-c19-s1:hmtch:prop-omega1-en-reales}
Dans ZFC, il existe une injection
\[
 \iota:\omega_1\hookrightarrow\mathbb R.
\]
Elle peut être relevée vers un sous-système binaire de l'espace des histoires HMT.
\end{proposition}

\begin{proof}
Par choix, la surjection \(\operatorname{ot}\) admet une section
\(s:\omega_1\to\operatorname{WO}\). De manière équivalente, on fixe un bon
ordre de \(\operatorname{WO}\) et l'on prend dans chaque fibre le premier code.
L'application de Cantor
\[
 c(a):=\sum_{n\geq0}\frac{2a(n)}{3^{n+1}}
\]
est injective de \(2^{\mathbb N}\) dans \([0,1]\). Par conséquent,
\[
 \iota:=c\circ s
\]
est l'injection recherchée. Le plongement du ruban binaire dans le ruban
ternaire et l'inverse de Hensel \(\Psi^{-1}\) fournissent le relèvement vers
l'espace généalogique.
\end{proof}

Le choix de \(s\) n'est pas naturel par rapport aux raffinements et n'est pas
produit par un préfixe fini. Cette proposition montre que la flèche
\(\omega_1\hookrightarrow\mathbb R\) n'est pas le contenu de l'hypothèse
classique du continu. La direction inverse appartient à l'ordinal extensionnel
nu. Dans HMT, en revanche, la flèche naturelle déjà construite est
l'homéomorphisme \eqref{p12-83-c18-s2:hmtch:eq-homeomorfismo-two-way-real}, dont le domaine
conserve les deux directions temporelles.

\section{Identité cardinale du continu analytique HMT : \(\aleph_{1,\mathrm{HMT}}^{\mathrm{an}}=2^{\aleph_0}\) et comparaison avec l'hypothèse classique du continu}

Soit
\[
 Q_{\mathrm{HMT}}:=X_{\mathrm{raw}}/\!\sim_P.
\]
D'après le \cref{p12-83-c18-s2:hmtch:thm-cociente-real},
\(Q_{\mathrm{HMT}}\cong\mathbb R\).

\begin{theorem}[Équivalences ordinales du quotient HMT]
\label{p12-83-c19-s1:hmtch:thm-equivalencias-ch}
Dans ZFC, sont équivalentes :
\begin{enumerate}
 \item \(2^{\aleph_0}=\aleph_1\);
 \item il existe une injection \(j:Q_{\mathrm{HMT}}\hookrightarrow\omega_1\) ;
 \item il existe une application \(\rho:X_{\mathrm{raw}}\to\omega_1\) telle que
 \[
  \rho(x)=\rho(y)
  \quad\Longleftrightarrow\quad
  P(x)=P(y);
 \]
 \item \(Q_{\mathrm{HMT}}\) admet un bon ordre non dénombrable dont les segments
 initiaux propres sont tous dénombrables.
\end{enumerate}
\end{theorem}

\begin{proof}
Si la première affirmation est satisfaite, une bijection
\(Q_{\mathrm{HMT}}\cong\mathbb R\cong\omega_1\) fournit la seconde et
transporte l'ordre de \(\omega_1\), ce qui donne la quatrième. La seconde donne la
troisième en prenant \(\rho=j\circ[P]\), où \([P]\) envoie chaque histoire sur sa
classe d'évaluation. Réciproquement, la condition biconditionnelle de la troisième
fait que \(\rho\) se factorise par une injection
\(Q_{\mathrm{HMT}}\hookrightarrow\omega_1\).

La seconde affirmation implique
\(2^{\aleph_0}\leq\aleph_1\). La
\cref{p12-83-c19-s1:hmtch:prop-omega1-en-reales} donne
\(\aleph_1\leq2^{\aleph_0}\) ; Cantor--Bernstein produit l'égalité.
Enfin, un bon ordre comme celui de la quatrième affirmation ne peut contenir un
élément à la position \(\omega_1\), car son segment initial serait non
dénombrable. Son type d'ordre est donc au plus \(\omega_1\), et puisque
le domaine est non dénombrable, il est exactement \(\omega_1\). Cela retrouve la
seconde affirmation.
\end{proof}

Le théorème localise la cible sans l'introduire comme prémisse. Une fonction
\(\rho\) possédant cette propriété ne serait pas un petit complément de la dichotomie
analytique : elle serait une formulation complète de l'hypothèse du
continu elle-même.

\section{Perte de mémoire généalogique sous réduction au codomaine ordinal}

Le foncteur d'oubli qui retire trit, miroir, quotients, route et mémoire envoie
chaque biordinal sur un support ordonné nu. Les résultats suivants
contrôlent ce foncteur d'oubli ; ils ne nient pas le morphisme HMT continu déjà
construit entre le quotient two-way et \(\mathbb R\).

\begin{theorem}[Limite des lecteurs vers \(\omega_1\) muni de la topologie d'ordre]
\label{p12-83-c19-s1:hmtch:thm-no-go-continuo}
Il n'existe pas d'application continue
\[
 \rho:X_{\mathrm{raw}}\longrightarrow\omega_1
\]
qui sépare exactement les fibres de \(P\). En particulier, après
application du foncteur d'oubli, aucun lecteur cylindrique vers l'ordinal
topologique nu ne peut démontrer la seconde condition du
\cref{p12-83-c19-s1:hmtch:thm-equivalencias-ch}.
\end{theorem}

\begin{proof}
Chaque carte \(\{m\}\times\mathbb Z_3^6\) est compacte. L'image continue d'un
compact dans \(\omega_1\), muni de la topologie d'ordre, est compacte et
est donc bornée par un certain ordinal dénombrable. Son image est dénombrable.
La réunion sur \(m\in\mathbb Z\) de ces images est une réunion dénombrable
d'ensembles dénombrables et demeure dénombrable. Une fonction qui séparerait toutes
les fibres de \(P\) aurait, en revanche, une image de cardinal
\(|\mathbb R|\), contradiction.
\end{proof}

\begin{theorem}[Absence de bon ordre analytique extensionnel]
\label{p12-83-c19-s1:hmtch:thm-no-go-analitico}
Il n'existe pas de bon ordre analytique, et donc pas non plus borélien, de
\(\mathbb R\).
\end{theorem}

\begin{proof}
Supposons que \(\prec\subseteq\mathbb R^2\) soit analytique et définisse un ordre
linéaire strict total. Alors
\[
 (x,y)\notin\prec
 \quad\Longleftrightarrow\quad
 x=y\ \text{ou}\ y\prec x.
\]
Le complémentaire de \(\prec\) est également analytique ; par séparation de Souslin,
\(\prec\) est borélien. Toute relation borélienne bien fondée sur un espace
polonais possède un rang borné par un ordinal dénombrable. Dans un bon ordre linéaire,
des éléments distincts ont des rangs distincts, de sorte que le domaine serait
dénombrable. Cela contredit la non-dénombrabilité de \(\mathbb R\).
\end{proof}

Les deux contrôles sont indépendants de la vérité ou de la fausseté de l'hypothèse classique du
continu. S'il existe une bijection ensembliste
\(\omega_1\leftrightarrow\mathbb R\), elle ne peut être continue pour la
topologie d'ordre de \(\omega_1\) ni analytique comme bon ordre. Cela ne rend pas
discontinu le morphisme HMT : cela montre que \(\omega_1\) nu est un autre
codomaine et que l'oubli de la direction temporelle détruit précisément la
structure qui assurait la continuité.

\section{Clôture ordinale du codomaine généalogique HMT}

Les contrôles précédents décrivent le foncteur qui oublie la généalogie. HMT
construit positivement sa clôture ordinale en la conservant. Soit
\(h\in2^{\mathbb N}\) le code engendré du langage, des tables finies,
des transitions et des lecteurs APP--TRIT--TPK. Le codomaine généalogique
canonique est
\[
 M_h:=L[h].
\]
Cette enveloppe conserve le générateur comme paramètre, admet le bon ordre
canonique de la hiérarchie constructible et est le codomaine produit par la
construction du continu dans \eqref{eq:frontal-doble-proyeccion}.

\begin{theorem}[Hypothèse classique du continu dans le codomaine généalogique HMT]
\label{p12-83-c19-s1:hmtch:thm-ch-relativa}
Le modèle \(M_h=L[h]\) satisfait
\[
 M_h\models 2^{\aleph_0}=\aleph_1.
\]
Il existe dans \(M_h\) une bijection définissable
\[
 b_h:\omega_1^{M_h}\xrightarrow{\ \sim\ }\mathbb R^{M_h}.
\]
Puisque \(M_h\) est le codomaine engendré, \(\mathbb R^{M_h}\) est la totalité de
la droite publiée par la construction HMT. Soit \(P_h\) l'évaluation HMT des
histoires appartenant à \(M_h\). Alors
\[
 \rho_h:=b_h^{-1}\circ P_h
\]
satisfait
\[
 \rho_h(x)=\rho_h(y)
 \quad\Longleftrightarrow\quad
 P_h(x)=P_h(y).
\]
\end{theorem}

\begin{proof}
Tout réel de \(L[h]\) apparaît à un certain niveau
\(L_\alpha[h]\) avec \(\alpha<\omega_1^{L[h]}\). Pour chaque ordinal
\(\alpha<\omega_1^{L[h]}\), le niveau \(L_\alpha[h]\) est dénombrable à l'intérieur
du modèle. Par conséquent,
\[
 |\mathbb R^{L[h]}|\leq\aleph_1^{L[h]}.
\]
La droite du modèle est non dénombrable, donc la minimalité de \(\aleph_1\)
donne l'inégalité inverse et l'on obtient l'égalité. Le bon ordre canonique
de \(L[h]\) sélectionne, parmi toutes les bijections dont l'existence vient
d'être prouvée, une première bijection définissable \(b_h\).

L'évaluation \(P_h\) est définissable à partir de \(h\), de sorte que sa valeur
appartient à \(\mathbb R^{M_h}\). La formule de \(\rho_h\) est bien typée.
Puisque \(b_h\) est bijective, deux histoires reçoivent le même ordinal si et seulement
si elles publient le même réel.
\end{proof}

\begin{remark}[Identification du codomaine engendré]
\label{p12-83-c19-s1:hmtch:rem-alcance-relativo}
Le \cref{p12-83-c19-s1:hmtch:thm-ch-relativa} construit effectivement la
flèche ordinale dans l'univers généalogiquement engendré. L'égalité
\(\mathbb R=\mathbb R^{L[h]}\) n'est pas postulée : dans le domaine du théorème,
elle exprime que \(L[h]\) est précisément le codomaine produit par le continu
HMT, dont la couverture archimédienne est prouvée dans les
\cref{p12-83-c18-s1:prop:agotamiento-recta-real-u024,p12-83-c18-s2:hmtch:thm-cociente-real}.
Par conséquent, \(\rho_h\) atteint tous les réels de la droite engendrée et
la bijection \(b_h\) démontre l'hypothèse classique du continu avec les
cardinaux ordinaires de \(M_h\). Une extension par forcing appartient à un autre
domaine après oubli ou remplacement du sélecteur généalogique \(h\) ; elle décrit
l'indépendance du langage nu de ZFC et ne modifie pas la conclusion du
système HMT.
\end{remark}

\section{Compatibilité entre la clôture cardinale HMT, la formulation ensembliste du continu et la mémoire généalogique bidirectionnelle}

Les résultats précédents articulent cinq strates d'une même clôture :
\begin{enumerate}
 \item le quotient généalogique reconstruit \(\mathbb R\) et a pour cardinal
       \(2^{\aleph_0}\) ;
 \item l'extension naturelle two-way et le biordinal généalogique sont
       homéomorphes, avec aller et retour continus ;
 \item le hamiltonien de l'horloge, le hamiltonien modulaire et le mode stable
       conservent calendrier, provenance et mémoire sans réduire le TPK au
       curseur visible ;
 \item toute publication HMT analytique est dénombrable ou possède le cardinal du
       continu, de sorte que
       \(\aleph_{1,\mathrm{HMT}}^{\rm an}=2^{\aleph_0}\) ;
 \item le codomaine constructible \(L[h]\) possède une énumération ordinale
       complète de tous les réels engendrés et satisfait l'hypothèse
       classique du continu avec ses cardinaux ordinaires.
\end{enumerate}
Les trois premières construisent l'objet HMT et sa dynamique ; la quatrième clôt la
puissance dans le domaine des publications ; la cinquième fournit sa coordonnée
ordinale sans changer d'objet. L'identité
\(\aleph_{1,\mathrm{HMT}}^{\rm an}=2^{\aleph_0}\), avec la bijection
\(b_h\), constitue la démonstration HMT de l'énoncé extensionnel sur la
droite engendrée. Le sélecteur \(h\) conserve la généalogie que le langage
nu de ZFC omet.

\begin{proposition}[Réfutateurs]
\label{p12-83-c19-s1:hmtch:prop-falsadores}
Chaque affirmation possède un contrôle négatif indépendant :
\begin{enumerate}
 \item un réel sans antécédent par \(P\) réfuterait la reconstruction du
       quotient ;
 \item un état intermédiaire non récupérable à partir des retours complets
       réfuterait la factorisation
       \eqref{p12-83-c17-s7:hmtch:eq-factorizacion-excepcional-retorno} ;
 \item la perte de la période primitive, de la reconstruction \(90/120\) ou
       du mode \(\varphi^{-2n}\) réfuterait l'un des trois lecteurs du
       registre temporel HMT typé ;
 \item une publication analytique non dénombrable sans copie de Cantor réfuterait
       la dichotomie parfaite ;
 \item un bon ordre borélien d'un espace polonais non dénombrable réfuterait le
       théorème d'impossibilité analytique ;
 \item une image continue non bornée d'une carte compacte dans \(\omega_1\)
       réfuterait le théorème d'impossibilité topologique ;
 \item l'existence, dans le codomaine engendré \(M_h\), d'un réel sans
       antécédent ordinal par \(b_h\), ou d'une branche publiée hors de
       \(L[h]\), réfuterait l'identification généalogique et la clôture classique
       du continu. Une extension obtenue après modification du sélecteur
       \(h\) n'appartient pas au domaine fixé et ne constitue pas ce réfutateur.
\end{enumerate}
\end{proposition}

La conclusion structurelle est précise. HMT construit un objet temporel
bidirectionnel, le publie archimédiennement comme \(\mathbb R\), conserve en
parallèle l'incidence exceptionnelle et démontre que sa première couche analytique
non dénombrable a pour cardinal exactement \(2^{\aleph_0}\). La bijection
\[
 \mathscr O_{\mathbb R,\leftrightarrow}^{\rm HMT}/\!
 \sim_{\mathfrak R}
 \xrightarrow{\ \overline{\mathfrak R}\ }\mathbb R
\]
est naturelle et continue. L'ordinal classique \(\omega_1^{M_h}\) est la coordonnée
de bon ordre de tous les réels du codomaine engendré. Le foncteur qui
oublie \(h\), le temps, le miroir et la mémoire explique l'indépendance ultérieure dans
ZFC nu ; il ne modifie ni la bijection construite ni la clôture classique
démontrée dans \(M_h\).
